\documentclass[11pt,a4paper]{article}

% -- Encodage & langue --------------------------------------------
\usepackage[utf8]{inputenc}
\usepackage[T1]{fontenc}
\usepackage[french]{babel}

% -- Mise en page ------------------------------------------------
\usepackage[top=2.5cm, bottom=2.5cm, left=2.8cm, right=2.8cm]{geometry}
\usepackage{parskip}
\setlength{\parskip}{6pt}

% -- Typographie -------------------------------------------------
% lmodern not available -- using default CM fonts
\usepackage[disable]{microtype}

% -- Couleurs ----------------------------------------------------
\usepackage[dvipsnames]{xcolor}
\definecolor{darkblue}{RGB}{26,58,92}
\definecolor{midblue}{RGB}{46,109,164}
\definecolor{lightblue}{RGB}{214,228,240}
\definecolor{accent}{RGB}{200,100,0}
\definecolor{darkgray}{RGB}{80,80,80}
\definecolor{lightgray}{RGB}{245,245,245}
\definecolor{warningbg}{RGB}{255,243,220}
\definecolor{warningborder}{RGB}{200,130,0}
\definecolor{okbg}{RGB}{232,245,238}
\definecolor{okborder}{RGB}{26,107,58}
\definecolor{nogobg}{RGB}{253,236,234}
\definecolor{nogoborder}{RGB}{192,57,43}
\definecolor{codebg}{RGB}{248,248,248}

% -- Sections ----------------------------------------------------
\usepackage{sectsty}
\usepackage{titlesec}
\titleformat{\section}{\large\bfseries\color{darkblue}}{}{0em}{}[\color{midblue}\rule{\linewidth}{0.5pt}]
\titleformat{\subsection}{\normalsize\bfseries\color{midblue}}{}{0em}{}
\titleformat{\subsubsection}{\normalsize\itshape\color{darkgray}}{}{0em}{}

% -- Tableaux ----------------------------------------------------
\usepackage{booktabs}
\usepackage{tabularx}
\usepackage{longtable}
\usepackage{array}
\usepackage{multirow}
\usepackage{colortbl}
\newcolumntype{L}[1]{>{\raggedright\arraybackslash}p{#1}}
\newcolumntype{C}[1]{>{\centering\arraybackslash}p{#1}}

% -- Boîtes colorées ---------------------------------------------
\usepackage{tcolorbox}
\tcbuselibrary{skins, breakable}

\newtcolorbox{infobox}[1][]{
  enhanced, breakable,
  colback=lightblue!40, colframe=midblue,
  boxrule=0.6pt, arc=4pt,
  left=8pt, right=8pt, top=6pt, bottom=6pt,
  fonttitle=\small\bfseries\color{darkblue},
  #1
}
\newtcolorbox{warningbox}[1][]{
  enhanced, breakable,
  colback=warningbg, colframe=warningborder,
  boxrule=0.6pt, arc=4pt,
  left=8pt, right=8pt, top=6pt, bottom=6pt,
  fonttitle=\small\bfseries\color{accent},
  #1
}
\newtcolorbox{okbox}[1][]{
  enhanced, breakable,
  colback=okbg, colframe=okborder,
  boxrule=0.6pt, arc=4pt,
  left=8pt, right=8pt, top=6pt, bottom=6pt,
  fonttitle=\small\bfseries\color{okborder},
  #1
}
\newtcolorbox{nogobox}[1][]{
  enhanced, breakable,
  colback=nogobg, colframe=nogoborder,
  boxrule=0.6pt, arc=4pt,
  left=8pt, right=8pt, top=6pt, bottom=6pt,
  fonttitle=\small\bfseries\color{nogoborder},
  #1
}

% -- Code --------------------------------------------------------
\usepackage{listings}
\lstset{
  backgroundcolor=\color{codebg},
  basicstyle=\ttfamily\small,
  keywordstyle=\color{midblue}\bfseries,
  commentstyle=\color{darkgray}\itshape,
  stringstyle=\color{accent},
  frame=single, framesep=4pt,
  frameround=tttt,
  rulecolor=\color{lightblue!80!black},
  breaklines=true,
  showstringspaces=false,
  xleftmargin=8pt, xrightmargin=4pt,
  aboveskip=8pt, belowskip=8pt,
}
\lstdefinelanguage{YAML}{
  keywords={true,false,null,yes,no},
  sensitive=false,
  comment=[l]{\#},
  morestring=[b]',
  morestring=[b]",
}

% -- Figures & graphiques ----------------------------------------
\usepackage{graphicx}
\usepackage{tikz}
\usetikzlibrary{shapes.geometric, arrows.meta, positioning, fit, backgrounds, calc}
\usepackage{pgfplots}
\pgfplotsset{compat=1.18}
\usepackage{float}
\usepackage{caption}
\captionsetup{font=small, labelfont=bf, labelsep=period}

% -- Maths --------------------------------------------------------
\usepackage{amsmath}
\usepackage{amssymb}

% -- Listes ------------------------------------------------------
\usepackage{enumitem}
\setlist[itemize]{topsep=2pt, itemsep=2pt, parsep=0pt}
\setlist[enumerate]{topsep=2pt, itemsep=2pt, parsep=0pt}

% -- Liens -------------------------------------------------------
\usepackage{hyperref}
\hypersetup{
  colorlinks=true,
  linkcolor=midblue,
  citecolor=midblue,
  urlcolor=midblue,
  pdftitle={Mini-SOC Low-Cost pour PME Togolaises},
  pdfauthor={Adrien},
}

% -- Bibliographie -----------------------------------------------
\usepackage{natbib}
\bibliographystyle{plainnat}

% -- En-têtes/pieds ----------------------------------------------
\usepackage{fancyhdr}
\pagestyle{fancy}
\fancyhf{}
\rhead{\small\color{darkgray}\thepage}
\lhead{\small\color{darkgray}Mini-SOC — EPL Lomé}
\renewcommand{\headrulewidth}{0.4pt}
\cfoot{\small\color{darkgray}Projet de remplacement de stage — L3 Informatique et Systèmes}

% -- Titre --------------------------------------------------------
\title{
  {\large\color{darkgray} Projet de remplacement de stage — L3 Informatique et Systèmes}\\[0.4cm]
  {\LARGE\bfseries\color{darkblue}
    Conception d'un Mini-SOC Low-Cost\\
    pour PME Togolaises}\\[0.3cm]
  {\large\color{midblue}
    Moteur de corrélation stateful en Python\\
    et évaluation comparative avec Wazuh}
}
\author{
  \textbf{Adrien}\\
  \small École Polytechnique de Lomé (EPL)\\
  \small L3 Informatique et Systèmes
}
\date{Mars 2026}

% ════════════════════════════════════════════════════════════════
\begin{document}
% ════════════════════════════════════════════════════════════════

\maketitle
\thispagestyle{empty}

% -- Résumé -------------------------------------------------------
\begin{abstract}
\noindent
Ce projet propose la conception et l'implémentation d'un \emph{mini-SOC open-source}
adapté aux contraintes des PME togolaises : budget limité, ressources matérielles
restreintes, et pénurie de compétences en cybersécurité.
Le système comprend un pipeline de collecte et normalisation de logs,
un moteur de corrélation stateful développé en Python avec règles externalisées en YAML,
et deux scénarios d'attaque simulés représentatifs des menaces locales.
La solution est rigoureusement évaluée et comparée à Wazuh, solution industrielle
open-source de référence, selon des métriques de détection (taux vrais/faux positifs),
de latence et de consommation de ressources.
Les livrables incluent un dépôt Git reproductible (IaC), un rapport analytique
et une démonstration technique en direct.
Ce projet s'inscrit dans le contexte de la Stratégie Nationale de Cybersécurité du Togo
2024--2028 et vise à démontrer la faisabilité d'une approche de détection
à coût quasi-nul pour les organisations de la sous-région.
\end{abstract}

\vspace{0.5cm}
\begin{infobox}[title={Note méthodologique importante}]
Les chiffres de performance présentés dans ce rapport (taux de détection, latence, RAM)
sont des \textbf{objectifs de conception}, pas des résultats mesurés.
Les mesures réelles seront collectées lors des phases de test (semaines~5--8)
et viendront remplacer ces valeurs cibles dans le rapport final.
\end{infobox}

\tableofcontents
\newpage

% ════════════════════════════════════════════════════════════════
\section{Introduction}
% ════════════════════════════════════════════════════════════════

\subsection{Contexte togolais}

\subsubsection{Données sourcées}

L'Union Internationale des Télécommunications (UIT) publie annuellement
l'indice mondial de cybersécurité (GCI), qui positionne le Togo parmi
les pays africains en progression sur les cinq piliers évalués
(cadre légal, capacités techniques, organisation institutionnelle,
renforcement des capacités, coopération internationale)~\citep{itu_gci2024}.
L'Agence Nationale de la Cybersécurité (ANCy), créée en 2020, constitue
le principal organe de coordination nationale en matière de sécurité numérique.

Sur le plan des coûts, les tarifs publics des solutions SIEM commerciales
permettent d'établir un constat objectif et vérifiable : une licence Splunk
Enterprise débute à plusieurs dizaines de milliers de dollars annuels,
et les solutions EDR professionnelles (CrowdStrike Falcon, SentinelOne)
facturent entre 50 et 150~USD par poste et par an.
Ces coûts sont structurellement incompatibles avec les capacités budgétaires
des petites et moyennes entreprises de la sous-région.

\subsubsection{Constats observationnels}

Les éléments suivants sont présentés comme des constats qualitatifs,
non comme des statistiques mesurées, en l'absence de sources académiques
publiées suffisamment granulaires sur le tissu des PME togolaises :

\begin{itemize}
  \item \textbf{Coût des solutions industrielles.} Les SIEM commerciaux
        (Splunk, IBM QRadar) et les EDR professionnels dépassent
        structurellement les capacités budgétaires des PME de la sous-région,
        indépendamment de toute statistique précise sur leur taux d'adoption.
  \item \textbf{Pénurie de compétences locales.} Les formations spécialisées
        en cybersécurité opérationnelle (Blue Team, analyse forensique,
        ingénierie SOC) sont peu développées en Afrique de l'Ouest
        francophone ; l'EPL Lomé constitue l'une des rares institutions
        proposant une filière informatique structurée au Togo.
  \item \textbf{Menaces peu couvertes localement.} Les vecteurs d'attaque
        les plus fréquents dans le contexte ouest-africain --- phishing via
        WhatsApp et SMS, fraude aux services Mobile Money (Flooz, T-Money),
        ransomwares ciblant des systèmes Windows non patchés --- sont peu
        représentés dans les bases de règles des solutions génériques
        internationales, conçues pour des contextes d'entreprise occidentaux.
\end{itemize}

\begin{infobox}[title={Positionnement méthodologique sur les données contextuelles}]
Ce rapport distingue explicitement deux niveaux de données :
\textbf{(1) données sourcées}, étayées par des publications officielles
citées en bibliographie ;
\textbf{(2) constats observationnels}, présentés comme tels
sans prétention statistique.
Cette distinction est intentionnelle : introduire des chiffres précis
sans source vérifiable dans un rapport académique constitue une faute
méthodologique. Si des rapports ANCy ou des études CEDEAO accessibles
sont identifiés avant la soutenance, ils remplaceront les constats
observationnels correspondants.
\end{infobox}

\subsection{Problématique}

\begin{center}
\begin{tikzpicture}
  \node[draw=midblue, fill=lightblue!30, rounded corners=6pt,
        text width=13cm, align=center, inner sep=10pt,
        font=\normalsize] (q) {
    \textbf{Question centrale}\\[4pt]
    Est-il possible de construire un système de détection d'intrusion stateful,
    capable d'identifier des chaînes d'attaques complexes (kill-chains),
    avec des contraintes matérielles compatibles avec les PME togolaises
    (\(\leq\)8\,Go RAM, coût logiciel nul), et dont les performances de détection
    sont comparables à celles de Wazuh sur les mêmes scénarios ?
  };
\end{tikzpicture}
\end{center}

\subsection{Hypothèses de travail}

\begin{enumerate}
  \item \textbf{H1 — Équivalence de détection.} Le moteur Python atteint un taux de
        vrais positifs (TPR) $\geq$\,90\,\% et un taux de faux positifs (FPR)
        $\leq$\,10\,\% sur les scénarios définis, sans différence statistiquement
        significative avec Wazuh sur ces mêmes scénarios.
  \item \textbf{H2 — Supériorité en ressources.} Le moteur Python consomme moins de
        4\,Go de RAM en charge et présente une latence de détection inférieure
        à celle de Wazuh sur le même matériel.
  \item \textbf{H3 — Adaptabilité locale.} Les règles de corrélation peuvent être
        écrites et modifiées par un utilisateur avec des compétences Python
        intermédiaires, sans expertise SIEM spécialisée.
\end{enumerate}

% ════════════════════════════════════════════════════════════════
\section{État de l'art}
% ════════════════════════════════════════════════════════════════

\subsection{Panorama des solutions existantes}

\begin{table}[H]
\centering
\caption{Comparaison des catégories de solutions de détection}
\small
\begin{tabularx}{\textwidth}{L{2.8cm} C{2cm} C{2cm} L{5cm}}
\toprule
\rowcolor{darkblue}
{\color{white}\textbf{Solution}} &
{\color{white}\textbf{Coût}} &
{\color{white}\textbf{RAM min.}} &
{\color{white}\textbf{Pertinence contexte PME Togo}} \\
\midrule
Splunk Enterprise & Très élevé & 16\,Go+ & Hors de portée financière \\
\midrule
IBM QRadar & Très élevé & 24\,Go+ & Hors de portée financière \\
\midrule
\rowcolor{lightgray}
Wazuh & Gratuit & 4\,Go+ & Viable techniquement, lourd en ressources \\
\midrule
Suricata (IDS réseau) & Gratuit & 1\,Go+ & Détection réseau uniquement, pas de corrélation \\
\midrule
\rowcolor{lightgray}
Antivirus basiques & Variable & Faible & Aucune corrélation, pas de détection réseau \\
\midrule
\textbf{Mini-SOC Python} & \textbf{Gratuit} & \textbf{$<$4\,Go} & \textbf{Cible de ce projet} \\
\bottomrule
\end{tabularx}
\end{table}

\begin{warningbox}[title={Pourquoi Norton est exclu du benchmark}]
Un antivirus grand public (Norton, McAfee, Avast) est fondamentalement différent
d'un moteur de corrélation stateful : il opère sur des signatures de fichiers et
des heuristiques comportementales sur un seul poste, sans capacité de corréler
des événements multi-sources dans le temps.
Comparer ces deux catégories d'outils n'aurait pas de sens technique.
Le benchmark retenu dans ce projet est \textbf{Wazuh uniquement},
qui appartient à la même catégorie fonctionnelle (SIEM/corrélation d'événements).
\end{warningbox}

\subsection{Wazuh — comparant de référence}

Wazuh est un SIEM open-source dérivé d'OSSEC, activement maintenu,
utilisé en production dans des milliers d'organisations.
Il comprend un agent déployé sur les hôtes supervisés, un manager
centralisant les événements, et un indexer (OpenSearch) pour le stockage
et la recherche.

Ses atouts pour ce projet : gratuité totale, présence de règles de détection
pré-écrites, documentation exhaustive, et usage croissant en Afrique
dans des structures à budget limité.
Sa limite principale : l'empreinte mémoire, notamment due à OpenSearch,
qui nécessite au minimum 4 à 8\,Go de RAM en configuration productive.

\subsection{Lacune identifiée}

Aucun outil de corrélation stateful n'existe à ce jour qui soit :
(a) entièrement open-source et gratuit,
(b) déployable sur moins de 4\,Go de RAM,
(c) dont les règles de détection soient lisibles et modifiables
    par un développeur Python sans formation SIEM,
(d) documenté avec des scénarios de menaces propres au contexte ouest-africain.

Ce projet adresse directement ce manque.

% ════════════════════════════════════════════════════════════════
\section{Architecture et méthodologie}
% ════════════════════════════════════════════════════════════════

\subsection{Vue d'ensemble du pipeline}

\begin{figure}[H]
\centering
\begin{tikzpicture}[
  node distance=0.6cm and 1.0cm,
  box/.style={
    draw=midblue, fill=lightblue!30, rounded corners=5pt,
    text width=2.6cm, align=center, minimum height=1.1cm,
    font=\small
  },
  arrow/.style={-Stealth, thick, color=midblue},
  label/.style={font=\scriptsize\itshape, color=darkgray},
]

\node[box] (logs)    {Logs système\\\& réseau};
\node[box, right=of logs]  (collect) {Collecte\\\small rsyslog};
\node[box, right=of collect] (norm)   {Normalisation\\JSON/ECS};
\node[box, right=of norm]   (engine) {Moteur Python\\stateful};
\node[box, right=of engine] (alert)  {Alertes\\\& réponse};

\draw[arrow] (logs)    -- (collect);
\draw[arrow] (collect) -- (norm);
\draw[arrow] (norm)    -- (engine);
\draw[arrow] (engine)  -- (alert);

\node[label, below=0.25cm of logs]    {Sources};
\node[label, below=0.25cm of collect] {Agrégation};
\node[label, below=0.25cm of norm]    {Standardisation};
\node[label, below=0.25cm of engine]  {Corrélation};
\node[label, below=0.25cm of alert]   {Sortie};

% Branche benchmark (en dessous)
\node[box, fill=warningbg!60, draw=warningborder,
      below=1.5cm of engine, text width=2.8cm] (wazuh)
      {Wazuh\\Manager};

\draw[-Stealth, thick, color=warningborder, dashed]
  (norm.south) .. controls +(0,-0.9) and +(-1.2,0) .. (wazuh.west)
  node[midway, label, left=2pt] {mêmes logs};

\draw[-Stealth, thick, color=warningborder, dashed]
  (wazuh.east) .. controls +(0.8,0) and +(0,-0.9) .. (alert.south)
  node[midway, label, right=2pt] {benchmark};

\end{tikzpicture}
\caption{Pipeline Mini-SOC et position du benchmark Wazuh}
\label{fig:pipeline}
\end{figure}

\subsection{Laboratoire virtualisé (Infrastructure as Code)}

\subsubsection{Topologie réseau}

Le laboratoire est entièrement isolé de l'Internet (réseaux VirtualBox
en mode \texttt{Host-Only} et \texttt{Internal Network}).
La topologie retenue est la suivante :

\begin{table}[H]
\centering
\caption{Composants du laboratoire}
\small
\begin{tabularx}{\textwidth}{L{3.2cm} C{1.6cm} C{2.2cm} L{4.8cm}}
\toprule
\rowcolor{darkblue}
{\color{white}\textbf{Composant}} &
{\color{white}\textbf{Type}} &
{\color{white}\textbf{IP}} &
{\color{white}\textbf{Rôle}} \\
\midrule
Firewall / Routeur & VM Alpine & 192.168.100.2 & Routage, segmentation, filtrage \\
\midrule
\rowcolor{lightgray}
Serveur collecte + moteur & Conteneur & 192.168.100.10 & Pipeline logs, moteur Python \\
\midrule
Cible SSH & Conteneur & 192.168.100.20 & Serveur SSH cible des attaques \\
\midrule
\rowcolor{lightgray}
Machine attaquante & VM Kali & 192.168.100.50 & Simulation des scénarios \\
\midrule
Wazuh Manager/Indexer & VM Debian & 192.168.100.30 & Benchmark (phases 8--10) \\
\bottomrule
\end{tabularx}
\end{table}

\subsubsection{Allocation mémoire}

Avec 16\,Go de RAM disponibles sur l'hôte Xubuntu, l'allocation est :

\begin{table}[H]
\centering
\caption{Budget RAM du laboratoire}
\small
\begin{tabularx}{\textwidth}{L{5cm} C{2cm} C{2cm}}
\toprule
\rowcolor{darkblue}
{\color{white}\textbf{Composant}} &
{\color{white}\textbf{RAM allouée}} &
{\color{white}\textbf{Phase active}} \\
\midrule
Firewall Alpine & 512\,Mo & 1--10 \\
\rowcolor{lightgray}
Serveur collecte + moteur & 1,5\,Go & 1--10 \\
Cible SSH & 512\,Mo & 1--10 \\
\rowcolor{lightgray}
Machine attaquante Kali & 1\,Go & 1--10 \\
Wazuh Manager & 2\,Go & 8--10 \\
\rowcolor{lightgray}
Wazuh Indexer (OpenSearch) & 2\,Go & 8--10 \\
\midrule
\textbf{Total phases 1--7} & \textbf{3,5\,Go} & — \\
\rowcolor{lightgray}
\textbf{Total phases 8--10} & \textbf{7,5\,Go} & — \\
\textbf{Marge système hôte} & \textbf{$\geq$8\,Go} & — \\
\bottomrule
\end{tabularx}
\end{table}

\subsection{Scénarios d'attaque retenus}

\begin{infobox}[title={Décision de scope — 2 scénarios uniquement}]
Le projet se limite à deux scénarios d'attaque définis dès la semaine~0
et figés jusqu'à la soutenance.
Cette décision est délibérée : deux scénarios maîtrisés, bien documentés,
avec des métriques précises, ont plus de valeur qu'une liste de scénarios
bâclés. Tout scénario supplémentaire éventuel sera traité en annexe
et clairement étiqueté « hors périmètre principal ».
\end{infobox}

\subsubsection{Scénario 1 — Brute-force SSH suivi d'une connexion réussie}

\textbf{Vecteur local :} Attaque automatisée sur des serveurs SSH exposés,
fréquente dans le contexte des petites structures togolaises dont les
services sont directement accessibles via IP publique sans VPN.

\textbf{Séquence d'événements :}
\begin{enumerate}
  \item Multiples tentatives de connexion SSH échouées depuis la même IP source
        ($\geq$\,10 tentatives en moins de 60\,secondes).
  \item Connexion SSH réussie depuis la même IP source dans les 120\,secondes
        suivant la dernière tentative échouée.
\end{enumerate}

\textbf{Règle de corrélation formelle :}
\[
  \exists\, t_0 : \bigl|\{e \mid e.\mathit{type}=\texttt{ssh\_fail},\,
  e.\mathit{src}=s,\, t_0 \leq e.t < t_0+60 \}\bigr| \geq 10
  \;\wedge\;
  \exists\, t_1 : e.\mathit{type}=\texttt{ssh\_ok},\, e.\mathit{src}=s,\,
  t_0 \leq t_1 \leq t_0+180
\]

\subsubsection{Scénario 2 — Reconnaissance réseau suivie d'une exfiltration}

\textbf{Vecteur local :} Attaque en deux temps visant des serveurs
de fichiers internes, représentative des intrusions observées dans les
PME de la sous-région (accès initial via credential faible, puis
exfiltration vers service cloud ou IP externe).

\textbf{Séquence d'événements :}
\begin{enumerate}
  \item Scan réseau large (SYN scan ou UDP scan) depuis une IP interne
        vers au moins 20 hôtes en moins de 30\,secondes.
  \item Transfert de données sortant vers une IP externe ($\geq$\,10\,Mo)
        dans les 300\,secondes suivant le scan.
\end{enumerate}

\textbf{Règle de corrélation formelle :}
\[
  \exists\, t_0 :
  \bigl|\{e \mid e.\mathit{type}=\texttt{scan},\, e.\mathit{src}=s,\,
  t_0 \leq e.t < t_0+30 \}\bigr| \geq 20
  \;\wedge\;
  \exists\, t_1 : e.\mathit{type}=\texttt{transfer\_out},\,
  e.\mathit{size} \geq 10^7,\, e.\mathit{src}=s,\,
  t_0 \leq t_1 \leq t_0+300
\]

\subsection{Architecture du moteur Python}

\subsubsection{Environnement technique et dépendances}

Le moteur est développé et testé sous \textbf{Python~3.12}
(version installée par défaut sur Ubuntu~24.04~LTS / Xubuntu~24.04,
environnement de développement de ce projet).
La version minimale requise est Python~3.10 (introduction des
\emph{union types} \texttt{X | Y} utilisés dans les annotations).
Toutes les dépendances sont fixées dans \texttt{requirements.txt}
à la racine du dépôt :

\begin{lstlisting}[language=bash, caption={requirements.txt -- dépendances exactes du moteur}]
# requirements.txt -- Python 3.12 (min. 3.10)
pyyaml>=6.0        # chargement et parsing des regles YAML
watchdog>=4.0      # surveillance inotify des fichiers de logs
jsonschema>=4.0    # validation du schema JSON des evenements
pytest>=8.0        # tests unitaires du moteur FSM
\end{lstlisting}

\begin{table}[H]
\centering
\caption{Stack technique du moteur Python -- rôle de chaque dépendance}
\small
\begin{tabularx}{\textwidth}{L{2.5cm} C{2.2cm} L{7.3cm}}
\toprule
\rowcolor{darkblue}
{\color{white}\textbf{Composant}} &
{\color{white}\textbf{Version}} &
{\color{white}\textbf{Rôle précis dans le moteur}} \\
\midrule
Python & 3.12 (min.\ 3.10) & Langage principal ; annotations de type \texttt{dict | None} requièrent 3.10+ \\
\rowcolor{lightgray}
PyYAML & $\geq$\,6.0 & Chargement des règles \texttt{*.yaml} depuis \texttt{engine/rules/} au démarrage \\
watchdog & $\geq$\,4.0 & Surveillance en temps réel du répertoire de logs via \texttt{inotify} (Linux) ; déclenche le parser à chaque nouvelle ligne \\
\rowcolor{lightgray}
jsonschema & $\geq$\,4.0 & Validation du schéma des événements normalisés avant injection dans la FSM \\
pytest & $\geq$\,8.0 & Tests unitaires : rejoue des séquences d'événements et vérifie les alertes générées \\
\bottomrule
\end{tabularx}
\end{table}

\subsubsection{Choix du backend d'états}

Deux options ont été évaluées pour le stockage des états de corrélation.

\paragraph{Option A --- Redis (écarté).}
Redis offre TTL natif, pub/sub et persistence RDB/AOF.
Il est écarté pour trois raisons : (1)~dépendance externe non justifiée
pour un moteur mono-serveur gérant deux scénarios ; (2)~overhead mémoire
de 50--200\,Mo inutile pour un lab ; (3)~états opaques sans \texttt{redis-cli},
ce qui alourdit les tests unitaires.
Cette décision est réversible : si le projet évolue vers plusieurs
sources de logs simultanées, Redis sera réévalué.

\paragraph{Option B --- Dictionnaire Python avec persistence JSON (retenu).}
Le moteur stocke les états actifs dans un dictionnaire Python en mémoire.
\textbf{Le risque principal est la perte d'état au redémarrage} : si le
processus plante, les attaques en cours de corrélation sont perdues.

La réponse retenue est une \emph{persistence périodique sur JSON}
avec reprise automatique au démarrage, implémentée dans la classe
\texttt{StateStore} décrite ci-dessous.
La fenêtre de perte maximale est bornée à 30\,secondes (paramétrable),
ce qui est acceptable au regard des fenêtres de détection de
60 à 300\,secondes définies dans les règles YAML.

\begin{lstlisting}[language=Python,
  caption={engine/state\_store.py -- Persistence JSON avec reprise au redémarrage}]
# engine/state_store.py  --  Python 3.12
import json, time, threading, pathlib

SNAPSHOT_PATH = pathlib.Path("engine/state_snapshot.json")
SNAPSHOT_INTERVAL = 30  # secondes entre chaque sauvegarde (parametrable)

class StateStore:
    """Dictionnaire d'etats avec persistence JSON periodique.

    Au demarrage : recharge le snapshot existant et purge
    les etats dont la fenetre temporelle est depassee.
    En fonctionnement : sauvegarde toutes les SNAPSHOT_INTERVAL secondes
    dans un thread daemon (n'empeche pas l'arret propre du processus).
    Perte maximale en cas de crash : SNAPSHOT_INTERVAL secondes d'etats.
    """

    def __init__(self):
        self._states: dict = {}
        self._lock = threading.Lock()
        self._load()          # reprise au demarrage
        self._start_loop()    # sauvegarde periodique en arriere-plan

    # -- Persistence --------------------------------------------------

    def _load(self) -> None:
        """Recharge le snapshot et purge les etats expires."""
        if not SNAPSHOT_PATH.exists():
            return
        with open(SNAPSHOT_PATH) as f:
            raw: dict = json.load(f)
        now = time.time()
        # Ne conserve que les etats dont la fenetre n'est pas depassee
        self._states = {
            k: v for k, v in raw.items()
            if now - v["last_seen"] < v["window_seconds"]
        }

    def _save(self) -> None:
        """Ecrit le snapshot sur disque (ecriture atomique via tmp)."""
        tmp = SNAPSHOT_PATH.with_suffix(".tmp")
        with self._lock:
            with open(tmp, "w") as f:
                json.dump(self._states, f, indent=2)
        tmp.replace(SNAPSHOT_PATH)  # atomique sur POSIX

    def _start_loop(self) -> None:
        """Lance le thread de sauvegarde periodique (daemon)."""
        def loop():
            while True:
                time.sleep(SNAPSHOT_INTERVAL)
                self._save()
        threading.Thread(target=loop, daemon=True).start()

    # -- Interface publique -------------------------------------------

    def get(self, key: str) -> dict | None:
        return self._states.get(key)

    def set(self, key: str, value: dict) -> None:
        with self._lock:
            self._states[key] = value

    def delete(self, key: str) -> None:
        with self._lock:
            self._states.pop(key, None)
\end{lstlisting}

\begin{infobox}[title={Ce que ce mécanisme garantit et ne garantit pas}]
\textbf{Garantit :}
\begin{itemize}[topsep=0pt]
  \item Reprise des états non expirés au redémarrage (crash ou mise à jour).
  \item Fenêtre de perte bornée à 30\,s (paramétrable via \texttt{SNAPSHOT\_INTERVAL}).
  \item Purge automatique des états dont la fenêtre temporelle est dépassée.
  \item Écriture atomique via fichier temporaire + \texttt{os.replace} :
        pas de snapshot corrompu en cas de crash pendant la sauvegarde.
\end{itemize}
\textbf{Ne garantit pas :}
\begin{itemize}[topsep=0pt]
  \item Scalabilité multi-processus (limitation connue, hors périmètre lab).
  \item Durabilité sub-seconde (Redis serait requis).
\end{itemize}
\textbf{Conclusion :} approche acceptable pour le périmètre mono-serveur de ce projet.
Décision documentée dans \texttt{docs/decisions.md} du dépôt Git.
\end{infobox}

\subsubsection{Structure des règles YAML}

Les règles de corrélation sont externalisées dans des fichiers YAML,
indépendants du code du moteur. Un exemple de règle pour le scénario~1 :

\begin{lstlisting}[language=YAML, caption={Règle YAML — Scénario 1 : brute-force SSH}]
# rules/ssh_bruteforce.yaml
id: SSH_BRUTEFORCE_001
name: "Brute-force SSH suivi d'une connexion reussie"
description: >
  Detecte une succession de tentatives SSH echouees
  suivie d'une connexion reussie depuis la meme source.
# -- Etape 1 : accumulation d'echecs --
trigger:
  event_type: ssh_auth_failed
  threshold: 10          # >= 10 echecs
  window_seconds: 60     # dans une fenetre de 60s
  group_by: src_ip       # par IP source
# -- Etape 2 : confirmation --
escalation:
  event_type: ssh_auth_success
  same_field: src_ip     # meme IP source
  within_seconds: 180    # dans les 3 minutes suivantes
# -- Sortie --
alert:
  severity: HIGH
  mitre_tactic: "TA0006"  # Credential Access
  message: "Kill-chain SSH brute-force detectee depuis {src_ip}"
\end{lstlisting}

\subsubsection{Architecture du moteur}

\begin{figure}[H]
\centering
\begin{tikzpicture}[
  mod/.style={
    draw=midblue, fill=lightblue!20, rounded corners=4pt,
    text width=3cm, align=center, minimum height=1cm, font=\small
  },
  container/.style={
    draw=darkgray, dashed, rounded corners=8pt,
    inner sep=10pt, fill=lightgray!30
  },
  arrow/.style={-Stealth, thick, midblue},
]

\node[mod] (parser)   {Parser de logs\\\small(JSON/Syslog)};
\node[mod, right=1.2cm of parser]  (norm)    {Normaliseur\\\small(ECS)};
\node[mod, right=1.2cm of norm]    (fsm)     {FSM stateful\\\small(dict Python)};
\node[mod, right=1.2cm of fsm]     (dispatch){Dispatcher\\\small(alertes)};

\node[mod, above=0.8cm of fsm, fill=warningbg!40,
      draw=warningborder, text width=2.6cm] (rules)
      {Règles YAML\\\small(externalisées)};

\draw[arrow] (parser)   -- (norm);
\draw[arrow] (norm)     -- (fsm);
\draw[arrow] (fsm)      -- (dispatch);
\draw[-Stealth, thick, warningborder] (rules) -- (fsm);

\begin{scope}[on background layer]
  \node[container, fit=(parser)(norm)(fsm)(dispatch)(rules),
        label={[font=\small\bfseries, color=darkgray]above:Moteur de corrélation}]
        (engine) {};
\end{scope}

\end{tikzpicture}
\caption{Modules internes du moteur de corrélation}
\label{fig:engine}
\end{figure}

% ════════════════════════════════════════════════════════════════
\section{Plan d'implémentation — 10 semaines}
% ════════════════════════════════════════════════════════════════

\begin{longtable}{C{1.5cm} L{3cm} L{7.5cm}}
\caption{Plan d'implémentation détaillé} \label{tab:plan} \\
\toprule
\rowcolor{darkblue}
{\color{white}\textbf{Sem.}} &
{\color{white}\textbf{Phase}} &
{\color{white}\textbf{Tâches principales}} \\
\midrule
\endfirsthead
\toprule
\rowcolor{darkblue}
{\color{white}\textbf{Sem.}} &
{\color{white}\textbf{Phase}} &
{\color{white}\textbf{Tâches principales}} \\
\midrule
\endhead
\midrule
\multicolumn{3}{r}{\small\itshape Suite page suivante\ldots}\\
\endfoot
\bottomrule
\endlastfoot

0 & Préparation & Dépôt Git + structure, schéma réseau, choix distributions, installation outils, décision VM/conteneur. \\
\rowcolor{lightgray}
1--2 & Topologie réseau IaC & Déploiement firewall Alpine, routage, segmentation. Docker Compose cibles + serveur collecte. Tests de connectivité. \\
3 & Pipeline collecte & rsyslog + journald → collecte centralisée → normalisation JSON/ECS. \textbf{Isolation du dataset d'évaluation (30\%).} \\
\rowcolor{lightgray}
4 & Moteur Python — base & Parser logs, machine à états (dict Python), fenêtres temporelles, chargement règles YAML. Tests unitaires. \\
5 & Scénario 1 + métriques & Simulation brute-force SSH. Mesure TPR/FPR. Ajustement seuils. \\
\rowcolor{lightgray}
6 & Scénario 2 + Wazuh & Simulation scan→exfiltration. \textbf{Déploiement Wazuh.} Calibration FP. Collecte métriques comparatives. \\
7 & Mini-SOC complet & Exposition alertes (fichier structuré/API). Mécanisme de réponse minimal (optionnel). \\
\rowcolor{lightgray}
8 & Benchmark Wazuh & Exécution des mêmes scénarios sur Wazuh. Mesure TPR, FPR, latence, RAM/CPU. Utilisation du dataset d'évaluation isolé. \\
9 & Analyse \& rapport & Comparaison chiffrée. Rédaction sections techniques. Visualisations. \\
\rowcolor{lightgray}
10 & Soutenance & Présentation + démonstration d'attaque/détection en direct (scénario~1 ou~2). \\

\end{longtable}

\begin{infobox}[title={Règle d'or — Dataset d'évaluation}]
Le dataset d'évaluation (30\,\% des logs générés) est isolé en semaine~3,
placé dans \texttt{datasets/eval/} et \textbf{jamais utilisé pendant le développement}.
Il sert exclusivement à l'évaluation finale du moteur Python ET de Wazuh.
Toute contamination de ce dataset invalide la comparaison et donc
la validité scientifique du benchmark.
\end{infobox}

% ════════════════════════════════════════════════════════════════
\section{Métriques d'évaluation}
% ════════════════════════════════════════════════════════════════

\subsection{Métriques de détection}

Les métriques suivantes seront calculées pour chaque scénario,
pour le moteur Python et pour Wazuh :

\begin{align}
  \text{TPR} &= \frac{VP}{VP + FN} \quad \text{(sensibilité)} \\
  \text{FPR} &= \frac{FP}{FP + VN} \quad \text{(taux de faux positifs)} \\
  \text{Précision} &= \frac{VP}{VP + FP}
\end{align}

où VP = vrai positif, FN = faux négatif, FP = faux positif, VN = vrai négatif.

\subsection{Métriques de performance}

\begin{itemize}
  \item \textbf{Latence de détection :} délai entre le dernier événement
        déclenchant et la génération de l'alerte (mesure en millisecondes,
        moyenne sur 20 répétitions par scénario).
  \item \textbf{RAM en charge :} consommation mémoire du processus Python
        ou du stack Wazuh pendant l'exécution des scénarios (via \texttt{/proc}).
  \item \textbf{CPU :} utilisation CPU moyenne sur la fenêtre de détection.
\end{itemize}

\subsection{Objectifs de conception (valeurs cibles)}

\begin{warningbox}[title={Rappel — Valeurs cibles, pas résultats}]
Le tableau ci-dessous présente des \textbf{objectifs de conception}.
Les valeurs réelles mesurées les remplaceront dans le rapport final.
Aucune de ces valeurs n'a été mesurée au moment de la rédaction de ce document.
\end{warningbox}

\begin{table}[H]
\centering
\caption{Objectifs de conception — à valider lors des tests}
\small
\begin{tabularx}{\textwidth}{L{3.5cm} C{2.2cm} C{2.2cm} C{2.2cm}}
\toprule
\rowcolor{darkblue}
{\color{white}\textbf{Métrique}} &
{\color{white}\textbf{Moteur Python}} &
{\color{white}\textbf{Wazuh}} &
{\color{white}\textbf{Source}} \\
\midrule
TPR (\%) & $\geq 90$ & À mesurer & Objectif H1 \\
\rowcolor{lightgray}
FPR (\%) & $\leq 10$ & À mesurer & Objectif H1 \\
Latence (ms) & $<3\,000$ & À mesurer & Objectif H2 \\
\rowcolor{lightgray}
RAM en charge & $<4$\,Go & À mesurer & Objectif H2 \\
\bottomrule
\end{tabularx}
\end{table}

% ════════════════════════════════════════════════════════════════
\section{Impact et perspectives locales}
% ════════════════════════════════════════════════════════════════

\subsection{Adéquation au contexte togolais}

Ce projet répond à une contrainte réelle du tissu économique togolais :
la nécessité de sécuriser des infrastructures modestes avec des moyens limités.
Les choix techniques reflètent cette contrainte à chaque niveau :

\begin{itemize}
  \item \textbf{Coût nul.} L'ensemble de la stack (Python, Docker,
        Alpine Linux, Wazuh) est open-source et ne nécessite aucune licence.
  \item \textbf{Empreinte mémoire maîtrisée.} Le moteur Python fonctionne
        sous 4\,Go, compatible avec du matériel de seconde main
        couramment disponible au Togo.
  \item \textbf{Règles adaptables.} Les règles YAML peuvent être
        rédigées et modifiées par un développeur avec des compétences
        Python, sans formation SIEM spécialisée.
  \item \textbf{Scénarios représentatifs.} Les deux scénarios retenus
        (brute-force SSH, exfiltration) correspondent aux vecteurs
        d'attaque les plus fréquents dans les PME de la sous-région
        disposant de services exposés avec des credentials faibles.
\end{itemize}

\subsection{Limites}

\begin{itemize}
  \item \textbf{Pas de scalabilité entreprise.} Le moteur Python en mémoire
        n'est pas conçu pour ingérer des milliers d'événements par seconde.
        Il cible explicitement les PME avec un trafic modéré.
  \item \textbf{Lab simulé vs production réelle.} Les métriques obtenues
        en laboratoire pourront différer du comportement en environnement
        de production réel (bruit de fond, anomalies inattendues).
  \item \textbf{Projet solo.} L'étendue du projet est limitée par le fait
        qu'un seul développeur en assure l'implémentation et les tests.
\end{itemize}

\subsection{Perspectives}

\begin{itemize}
  \item Collaboration avec l'ANCy (Agence Nationale de la Cybersécurité)
        pour validation et éventuelle diffusion.
  \item Déploiement pilote au sein de l'EPL Lomé sur
        l'infrastructure réseau interne.
  \item Extension à d'autres scénarios de menaces mobiles
        (fraude Mobile Money, phishing WhatsApp) dans une version~2.
  \item Contribution à un référentiel de règles de détection
        contextualisées pour l'Afrique de l'Ouest.
\end{itemize}

% ════════════════════════════════════════════════════════════════
\section{Conclusion}
% ════════════════════════════════════════════════════════════════

Ce projet démontre la faisabilité d'une approche de détection d'intrusion
stateful, low-cost et adaptée au contexte des PME togolaises.
En construisant un moteur de corrélation Python avec des règles externalisées
et en le benchmarquant rigoureusement contre Wazuh, il produit
à la fois un livrable technique concret et une analyse comparative
méthodologiquement solide.

Les compétences développées — architecture pipeline de données,
programmation Python orientée sécurité, infrastructure as code,
conduite de tests comparatifs — sont directement transférables
à des rôles d'ingénieur SOC, d'analyste sécurité ou de consultant
en cybersécurité, aussi bien au Togo que dans la sous-région
ou à l'international.

% ════════════════════════════════════════════════════════════════
\nocite{*}
\bibliography{refs}
% ════════════════════════════════════════════════════════════════

\newpage
% ════════════════════════════════════════════════════════════════
\appendix
\section{Structure du dépôt Git}
% ════════════════════════════════════════════════════════════════

\begin{lstlisting}[language=bash, caption={Structure du dépôt mini-soc-togo}]
mini-soc-togo/
|---- infra/
|   |---- docker-compose.yml      # Stack collecte + cibles
|   \---- firewall/               # Config Alpine iptables
|---- collector/
|   |---- rsyslog.conf            # Config collecte
|   \---- normalizer.py           # Pipeline JSON/ECS
|---- engine/
|   |---- core.py                 # Moteur FSM principal
|   |---- loader.py               # Chargement regles YAML
|   |---- rules/
|   |   |---- ssh_bruteforce.yaml # Scenario 1
|   |   \---- scan_exfil.yaml     # Scenario 2
|   \---- tests/
|       |---- test_ssh.py         # Tests unitaires S1
|       \---- test_scan.py        # Tests unitaires S2
|---- scenarios/
|   |---- run_scenario1.sh        # Script attaque S1
|   \---- run_scenario2.sh        # Script attaque S2
|---- datasets/
|   |---- dev/                    # 70% - developpement
|   \---- eval/                   # 30% - NE PAS MODIFIER
|---- wazuh/
|   \---- custom-rules.xml        # Regles Wazuh equivalentes
|---- results/
|   \---- metrics.csv             # Resultats benchmark
|---- report/
|   |---- main.tex                # Ce document
|   \---- refs.bib                # Bibliographie
\---- docs/
    |---- decisions.md            # Journal decisions archi
    \---- network-topology.md     # Topologie reseau
\end{lstlisting}

% ════════════════════════════════════════════════════════════════
\section{Règles de corrélation — Scénario 2}
% ════════════════════════════════════════════════════════════════

\begin{lstlisting}[language=YAML, caption={Règle YAML — Scénario 2 : scan → exfiltration}]
# rules/scan_exfil.yaml
id: SCAN_EXFIL_001
name: "Reconnaissance reseau suivie d'une exfiltration"
description: >
  Detecte un scan reseau large depuis un hote interne
  suivi d'un transfert de donnees sortant vers une IP externe.
trigger:
  event_type: network_scan
  threshold: 20            # >= 20 hotes scanes
  window_seconds: 30       # en moins de 30 secondes
  group_by: src_ip
escalation:
  event_type: data_transfer_outbound
  same_field: src_ip
  min_bytes: 10000000      # >= 10 Mo
  within_seconds: 300      # dans les 5 minutes suivantes
alert:
  severity: CRITICAL
  mitre_tactic: "TA0010"  # Exfiltration
  message: "Kill-chain scan+exfiltration depuis {src_ip}"
\end{lstlisting}

% ════════════════════════════════════════════════════════════════
\section{Bibliographie — Sources à compléter}
% ════════════════════════════════════════════════════════════════

\begin{warningbox}[title={Sources à vérifier et compléter avant le rapport final}]
Les références ci-dessous sont des pointeurs à valider.
Chaque entrée doit être vérifiée et remplacée par une source accessible
et citable (rapport officiel, publication académique, site institutionnel).
\end{warningbox}

Sources à investiguer :
\begin{itemize}
  \item Rapports annuels de l'ANCy (Agence Nationale de la Cybersécurité du Togo)
  \item Rapports ITU sur l'index de cybersécurité mondial (GCI) — section Afrique
  \item Études CEDEAO sur la cybersécurité des PME en Afrique de l'Ouest
  \item Documentation officielle Wazuh : \url{https://documentation.wazuh.com}
  \item MITRE ATT\&CK Enterprise : \url{https://attack.mitre.org}
  \item RFC~5424 (Syslog protocol), RFC~8259 (JSON)
\end{itemize}

\end{document}